\documentclass[10pt,a4paper]{article}
\usepackage{amsmath,amssymb,amsthm,graphicx}
\usepackage{hyperref}
\usepackage{latexsym}
\usepackage{amssymb}
\usepackage{amsthm}
\usepackage{ulem}
\usepackage{enumitem}
\setlist{itemsep=0.4pt,topsep=0pt}
\parindent 0pt
\parindent 0pt
\parskip 12pt
\newcommand{\HCF}{\mathrm{hcf}}
\newcommand{\ZZ}{\mathbb{Z}}
\newcommand{\rk}{\mathrm{rk}}
\DeclareMathOperator{\len}{length}
\DeclareMathOperator{\bl}{b}
\DeclareMathOperator{\cl}{cl}
\DeclareMathOperator{\di}{d}
\DeclareMathOperator{\Cl}{Cl}
\DeclareMathOperator{\acl}{acl}
\DeclareMathOperator{\tp}{tp}
\DeclareMathOperator{\Th}{Th}
\newcommand{\Cor}{\mathcal{C}_\mathrm{or}}
\newcommand{\Cun}{\mathcal{C}_\mathrm{un}}
\newcommand{\mons}{\mathbb{M}}
\newcommand{\strong}{\mathrm{strong}}
\newcommand{\axioms}{\Delta_{\mathrm{FSTS}}}
\newtheoremstyle{note}% hnamei
{8pt}% hSpace abovei
{14pt}% hSpace belowi
{\normalfont}% hBody fonti
{}% hIndent amounti1
{\bfseries}% hTheorem head fonti
{ }% hPunctuation after theorem headi
{.5em}% hSpace after theorem headi2
{}% hTheorem head spec (can be left empty, meaning ‘normal’)i
\theoremstyle{note}
\newtheorem{thm}{Theorem}[section]
%\newtheorem{theorem}{Theorem}
\newtheorem{lemma}[thm]{Lemma} 
\newtheorem{df}[thm]{Definition}
\newtheorem{rmk}[thm]{Remark} 
\newtheorem{prop}[thm]{Proposition}
\newtheorem{fact}[thm]{Fact}
\newtheorem{cor}[thm]{Corollary}
\newtheorem{conj}[thm]{Conjecture}
\newtheorem{ex}[thm]{Example}
\newtheorem{qu}[thm]{Question}
\newtheorem*{notation}{Notation} 
\title{There is a countable unconfined Steiner triple system that is not free}
\author{Silvia Barbina, Enrique Casanovas}
\date{arxiv:2411.13723v3; Creative Commons Attribution 4.0 International}
\begin{document}
\maketitle
\noindent\textit{Editorial note.} The counted statement is the author's environment at source0.tex lines 1097-1098 and its proof environment opens at line 1099; the pairing is by adjacency and the retained proof range is the whole environment.\par\medskip
\noindent\textit{Editorial note.} The preamble is the author's own package lines, macro block and theorem-environment declarations (source lines 4-65); the author's optioned amsmath line and the fleqn-dependent layout line are not retained.\par\medskip
\noindent\textit{Editorial note.} The source carries no label-number dump and no counter values are imposed, so the printed numbering is the wrapper's default; identity is the author label, or the pinned source position where the author gave none.\par\medskip
\noindent\textit{Editorial note.} Metadata provenance: the source record of the pinned copy and the licence evidence file are both present; the licence is CC BY 4.0 with status supported, and the card's source\_version matches the value recorded in that evidence file.\par\medskip
\noindent\textit{Editorial note.} r002 note: the retained author proof cites the key barcas2, which the article bibliography bundled with the source carries; r002 bundles that bibliography file and changes nothing else.\par\medskip
\noindent\textit{Editorial scope.} One original result of the pinned author manuscript, counted once: a countable unconfined system that is not free (unlabelled proposition at source0.tex lines 1097-1098).\par\medskip
\noindent\textit{Allowed standard prerequisites (editorial).}\begin{itemize}
\item The free system on three generators. F\_3 is freely generated by three elements a\_1, a\_2, a\_3.. It is the starting point of the construction.
\item The iterated union. N\_omega is the union of a countable chain of copies of the free system, each elementary in the next.. It is the system shown to be unconfined but not free.
\item Freely generated systems. A system is free when it is freely generated by some set, in particular it is generated by that set.. It is what the constructed system fails to be.
\end{itemize}\medskip
	\begin{prop}\label{April_11_2024_6}
		Let $M,N$ be models of the theory of free STSs. If $M\leq N$, then  $M\preccurlyeq N$.
	\end{prop}

	\begin{prop}\label{April_11_2024_7} If $3\leq n<m$, there is no elementary embedding of $F_m$, the free STS with $m$ generators, into $F_n$, the free STS of $n$ generators, despite the elementary equivalence of $F_m$ and $F_n$.
	\end{prop}

	\begin{prop}   There is a (countable) unconfined STS which is not free.
	\end{prop}

	\begin{proof} Let $N=F_3$ be  the free STS generated by three elements  $a_1,a_2,a_3$, and let \\ $b = (a_1 \cdot a_2) \cdot (a_1 \cdot a_3)$. As shown in~\cite{barcas2},  Example~6.9, the homomorphism $f:N\rightarrow N$ determined by $f(a_1)=b$, $f(a_2)=a_2$ and $f(a_3)=a_3$ is an embedding but it is not surjective. Hence, there is some proper substructure of $N$ isomorphic to $N$. So we can define an ascending chain $(N_n\mid n<\omega)$ of STSs such that $N_n\cong N$ and $N_n\neq N_{n+1}$. Let $N_\omega=\bigcup_{n<\omega}N_n$. %Then $N_\omega$ is unconfined 
 Since $N_\omega$ is a union of a countable chain of unconfined STSs and unconfinedness fails on a finite subset, we have that $N_\omega$ is unconfined. Moreover, it is not finitely generated. Since it is countable, if it is free then it has a free base $B$ of $\omega$ elements and  it is (isomorphic to)  the generic. There is an embedding $f:F_4\rightarrow N_\omega$ such that  $f(F_4)\leq N_\omega$.  Choose $n$ such that the image in $f$ of  the base of $F_4$  is contained in $N_n$.  Then $f(F_4)\leq N_n$  and by Proposition~\ref{April_11_2024_6},  $f:F_4\rightarrow N_n$ is elementary. Since  $N_n\cong F_3$, this contradicts Proposition~\ref{April_11_2024_7}.
		Hence, $N_\omega$ is not free.
	\end{proof}

\begin{thebibliography}{9}
\bibitem{barcas2}
S.~Barbina and E.~Casanovas.
\newblock Endomorphisms of free {S}teiner quasigroups, 2024.
\newblock{\url{arXiv:2412.06055}}.

\bibitem{casanovas}
E.~Casanovas.
\newblock Compactly expandable models and stability.
\newblock {\em The Journal of Symbolic Logic}, 60:673--683, 1995.
\bibitem{ganter1}
B.~Ganter and H.~Werner.
\newblock Equational classes of {S}teiner systems.
\newblock {\em Algebra Universalis}, 5:125--140, 1975.

\bibitem{ganter2}
B.~Ganter and H.~Werner.
\newblock Co-ordinatizing {S}teiner systems.
\newblock In C.C. Linder and A.~Rosa, editors, {\em Topics on {S}teiner
  systems}, pages 3--24. North Holland, 1980.
\end{thebibliography}
\end{document}
